\section{Ameaças à validade}

\paragraph{Construto.} Nenhuma das métricas é observada diretamente. Tratar uma release como deploy só vale se ela chega a quem usa o projeto: o consenso da amostra-ouro foi que sim em 59 dos 60 repositórios, mas com kappa de 0,19 e cobrindo só 60 dos 300 (no \texttt{supabase}, as ``releases'' são boletins mensais). O CFR (a) mede falha de CI, e não de produção; o CFR (b) depende de uma heurística com F1 de 0,644, abaixo da meta de 0,70, com erros nos dois sentidos. O lead time usa a data de autoria do commit, que pode anteceder a sua entrada no ramo (\textit{rebase}, \textit{cherry-pick}), e 6 releases com a tag apagada ou reescrita ficaram fora. A recuperação mede a volta do CI, e não do serviço, e os 0,9\% de episódios censurados não entram na mediana. A RQ 07 dimensiona essa ameaça: a definição muda a categoria de 21,7\% a 30,1\% dos repositórios.

\paragraph{Interna.} A idade tem o maior efeito da RQ 06 (lead time, $\varepsilon^2 = 0{,}365$) e pode se sobrepor a outros fatores, o que a análise por fator isolado não separa. As correlações da RQ 05 não implicam causalidade, e o $\rho$ de 0,501 com o CFR (b) é em parte efeito da definição. As regras da heurística foram escolhidas a partir dos erros nos mesmos 300 itens usados para avaliá-la, o que torna o F1 otimista, e os autores foram também os avaliadores; como mitigação, o protocolo foi escrito antes do sorteio e a rotulagem foi independente e sem ver a saída da heurística.

\paragraph{Externa.} A amostra reúne os 300 repositórios mais populares que passam no funil, e não é um sorteio; as exigências de GitHub Actions, de 5 releases e de 50 execuções favorecem projetos que publicam releases no GitHub. Os resultados não se estendem a projetos pequenos, corporativos ou de outras plataformas, nem necessariamente a outros períodos, e 14 organizações têm mais de um repositório (até 6), o que reduz a independência das observações.

\paragraph{Conclusão.} Só 4 dos 20 testes da RQ 06 permanecem significativos após a correção de Holm, e os do tipo do projeto, com 58 repositórios (o grupo \textit{outro}, com 2, ficou abaixo do mínimo de 5), têm pouco poder. Os quartis são um agrupamento arbitrário, e os testes não paramétricos não ajustam um fator pelos demais. A correlação de Spearman tem muitos empates (11 repositórios com CFR (a) igual a zero e 61 com CFR (b) igual a zero), e o kappa de Fleiss é sensível à prevalência, por isso o de entregas reais é relatado com o acordo bruto (88\%). Por fim, valores próximos de um corte podem mudar de categoria por diferenças pequenas.
